\chapter{Results and scientific status}\label{ch:results}
\section{Level A: verified dispersion}
For the recorded representative inputs $d=\SI{1}{mm}$,
$\rho=\SI{2700}{kg.m^{-3}}$, $c_L=\SI{6320}{m.s^{-1}}$, and
$c_T=\SI{3130}{m.s^{-1}}$, the characteristic residual is below
$2.44\times10^{-9}$. Low-$k$ median relative errors are 0.011\% for S0 and
0.150\% for A0. The isolated first symmetric minimum is
$fd=\SI{2.8568}{MHz.mm}$ at $kh=0.7989$. This is consistent with the ZGV
frequency scale reported for aluminum plates by Prada et al. \cite{prada2005};
all material differences must be retained in any quantitative comparison.

\begin{figure}[htbp]\centering
\includegraphics[width=.92\linewidth]{dispersion.png}
\caption{Verified numerical Rayleigh--Lamb roots. Higher-order curves are broken
where frequency continuation is discontinuous; eigenfield-overlap tracking is
required before those segments are assigned as complete physical branches.}
\label{fig:dispersion}
\end{figure}

\section{Level B: physical laser coupling}
The governing overlap is established analytically in
\cref{eq:physical-overlap}. Numerical overlap curves, beam-width and absorption-
depth sweeps, and dimensional amplitudes are \emph{pending}. They require
normalized eigenfields plus the optical and thermal inputs listed in
\texttt{SCIENTIFIC\_STATUS.md}. No scalar unit-overlap output is reported as a
physical result.

\section{Level C: causal steady-state formation}
The exact modal decomposition and robust measurement definition are established
in \cref{ch:causal}. Existing scalar-overlap calculations are archived under
\texttt{results/preliminary}; their earlier ``modal purity'', 50-cycle settling
example, and parameter sweep are illustrative regression data only. They use
prescribed damping, equal scalar overlaps, and an uncalibrated spectral window.
They cannot validate experimental $N_{95}$ or selectivity. Physical results near
ordinary, cutoff, and ZGV points remain pending a detector model, calibrated
attenuation, and the level-B overlaps.
